\documentclass{article}
\usepackage[UTF8]{ctex}
\usepackage{amsmath}
\usepackage{amssymb}
\usepackage{algorithm}
\usepackage{algpseudocode}
\usepackage{graphicx}
\usepackage{geometry}
\usepackage{float}
\usepackage{hyperref}
\usepackage{booktabs}
\geometry{a4paper, margin=1in}
\title{Colormap 批量生成系统}
\begin{document}
\maketitle

\section{引言}
为心理物理学实验提供刺激材料，需要批量生成具有系统性参数变化且满足感知质量约束的色彩映射（Colormap）。本系统在 CIE HCL 色彩空间中工作，以色相（Hue）、色度（Chroma）和亮度（Luminance）的控制点曲线为输入，通过网格搜索 + 贪心优化 + 局部精炼的三阶段流水线生成候选色彩映射，再用基于 CIEDE2000 的双窗口均匀性检验进行质量筛选，最终得到 48 条高质量色彩映射（3 色相 $\times$ 4 色度模式 $\times$ 4 亮度模式）。

整体流程：
\begin{enumerate}
    \item \textbf{参数空间定义}：预定义色相目标、色度/亮度控制点模式
    \item \textbf{候选生成}：在 HCL 空间中网格枚举参数组合，合成 256 色的连续色彩映射
    \item \textbf{可行性过滤}：丢弃包含超出 sRGB 色域颜色的候选；通过 CIEDE2000 可辨别性检验
    \item \textbf{质量检验}：双窗口（小窗 / 大窗）CIEDE2000 均匀性检验
    \item \textbf{扰动优化}：贪心搜索最优扰动偏移 $(\delta_C, \delta_L)$，配合局部精炼寻找满足质量约束的最小偏移解
    \item \textbf{最优选择}：从多次生成结果中按最小扰动距离选出每个参数组合的最优个体
\end{enumerate}


\section{设计变量与参数空间}

\subsection{色彩空间：CIE HCL}
所有色彩映射在 CIE HCL（Hue--Chroma--Lightness）空间中定义，该空间是 CIELAB 的极坐标表示，具有更直观的感知解耦（色相、饱和度、明度相互正交）。HCL 坐标 $(H, C, L)$ 与 sRGB 之间通过 D3.js 库进行转换。

\subsection{色相目标}
预定义三个色相目标值：
\begin{equation}
H_{target} \in \{100°,\; 200°,\; 300°\}
\end{equation}
这三个值大致对应黄绿色、青蓝色和紫红色区域，覆盖了色相环的主要感知区间。色相范围 $\Delta H = H_{target}$ 定义了从起始色相到终止色相的跨度。

\subsection{色度与亮度控制点模式}
系统定义了四种控制点模式，分别应用于色度通道 $C(t)$ 和亮度通道 $L(t)$，其中 $t \in [0, 1]$ 为沿色彩映射的参数化位置：

\begin{table}[H]
\centering
\begin{tabular}{llll}
\toprule
模式名称 & 控制点序列 & 波形形状 & 语义 \\
\midrule
Constant  & $[v, v]$ & 水平直线 & 均匀 \\
Linear    & $[v_{lo}, v_{hi}]$ & 单调递增 & 渐变 \\
Diverging & $[v_{lo}, v_{hi}, v_{lo}]$ & 先升后降 & 中心聚焦 \\
Thermal   & $[v_{lo}, v_{hi}, v_{lo}, v_{hi}, v_{lo}]$ & 双峰振荡 & 热力图 \\
\bottomrule
\end{tabular}
\caption{控制点模式定义}
\end{table}

色度通道的具体控制点值为：
\begin{align}
C_{\text{constant}} &= [70, 70], \quad C_{\text{linear}} = [20, 120] \nonumber \\
C_{\text{diverging}} &= [20, 120, 20], \quad C_{\text{thermal}} = [20, 120, 20, 120, 20]
\end{align}

亮度通道的具体控制点值为：
\begin{align}
L_{\text{constant}} &= [55, 55], \quad L_{\text{linear}} = [20, 90] \nonumber \\
L_{\text{diverging}} &= [20, 90, 20], \quad L_{\text{thermal}} = [20, 90, 20, 90, 20]
\end{align}

\subsection{参数组合}
总计 $3 \times 4 \times 4 = 48$ 种参数组合，每种组合需生成一条满足质量约束的色彩映射。

\subsection{控制点调整机制}
控制点模式定义的是"理想"曲线形状。实际生成时，需要通过扰动偏移 $(\delta_C, \delta_L)$ 缩小控制点的值域范围，以增加在 sRGB 色域内找到可行解的概率。调整规则为：
\begin{equation}
C'_{max} = \max(C_{min},\; C_{max} - \delta_C), \quad L'_{max} = \max(L_{min},\; L_{max} - \delta_L)
\end{equation}
即仅将最大值控制点向下压缩 $\delta$ 个单位，保持最小值和曲线形状不变（Constant 模式不受扰动影响）。


\section{色彩映射合成算法}

\subsection{色彩通道插值}
给定控制点序列 $\mathbf{p} = [p_0, p_1, \dots, p_n]$，沿 $t \in [0, 1]$ 的分段线性插值为：
\begin{equation}
f(t; \mathbf{p}) = p_{\lfloor nt \rfloor} \cdot (1 - t_{local}) + p_{\lfloor nt \rfloor + 1} \cdot t_{local}, \quad t_{local} = nt - \lfloor nt \rfloor
\end{equation}

\subsection{单条色彩映射合成}
对于给定的色相起点 $H_0$、色度基准 $C_{base}$、亮度基准 $L_{base}$，系统采样 $N = 256$ 个色彩点：

对第 $i$ 个采样点（$i = 0, 1, \dots, 255$），参数化位置 $t = i / 255$：
\begin{align}
H_i &= H_0 + \Delta H \cdot t \pmod{360} \\
C_i &= C_{base} + \Delta C \cdot \hat{f}_C(t) \\
L_i &= L_{base} + \Delta L \cdot \hat{f}_L(t)
\end{align}
其中 $\hat{f}_C(t)$ 和 $\hat{f}_L(t)$ 是将控制点插值归一化到 $[0, 1]$ 后的函数：
\begin{equation}
\hat{f}(t) = \frac{f(t; \mathbf{p}) - \min(\mathbf{p})}{\max(\mathbf{p}) - \min(\mathbf{p})}
\end{equation}
$\Delta C = C_{max}' - C_{min}'$ 和 $\Delta L = L_{max}' - L_{min}'$ 为调整后的通道值域范围。

对于 Constant 模式（$\max = \min$），直接使用固定值而非插值。

每个 HCL 点 $(H_i, C_i, L_i)$ 通过色彩空间转换映射为 sRGB 值。若任意一个采样点超出 sRGB 可显示色域（$R, G, B \notin [0, 255]$），则\textbf{整条色彩映射被丢弃}。

\subsection{候选生成：网格搜索}
系统在三维参数空间 $(H_0, C_{base}, L_{base})$ 上进行网格枚举：

\begin{algorithm}[H]
\caption{色彩映射候选网格搜索}
\begin{algorithmic}[1]
\State \textbf{Input:} 目标色相范围 $\Delta H$，调整后的色度/亮度控制点
\State 随机生成色相起点种子 $H_{seed} \sim \mathcal{U}(0°, 360°)$
\For{$H_{offset} = 0°, 5°, 10°, \dots, 360°$} \Comment{步长 $= 5°$}
    \For{$C_{base} = 20, 25, 30, \dots, 130 - \Delta C$} \Comment{步长 $= 5$}
        \For{$L_{base} = 20, 25, 30, \dots, 100 - \Delta L$} \Comment{步长 $= 5$}
            \State 合成 256 色的色彩映射
            \If{所有颜色在 sRGB 色域内 \textbf{AND} 通过可辨别性检验}
                \State 加入候选集合
            \EndIf
        \EndFor
    \EndFor
\EndFor
\State 按色度/亮度目标匹配误差排序，返回最佳候选
\end{algorithmic}
\end{algorithm}


\section{质量控制}

\subsection{CIEDE2000 色差度量}
系统采用 CIEDE2000 标准色差公式作为所有质量判断的基础度量：
\begin{equation}
\Delta E_{00} = \sqrt{
    \left(\frac{\Delta L'}{k_L S_L}\right)^2 +
    \left(\frac{\Delta C'}{k_C S_C}\right)^2 +
    \left(\frac{\Delta H'}{k_H S_H}\right)^2 +
    R_T \frac{\Delta C'}{k_C S_C} \frac{\Delta H'}{k_H S_H}
}
\end{equation}
其中 $S_L, S_C, S_H$ 为与平均亮度/色度相关的缩放因子，$R_T$ 为蓝色区域旋转项，$k_L = k_C = k_H = 1$。该公式相比简单的欧氏 $\Delta E_{ab}$ 更接近人类感知的色差判断。

\subsection{基础可辨别性检验}
在候选生成阶段，首先进行粗粒度的可辨别性过滤。从 256 色中均匀采样 10 个颜色，要求\textbf{任意两个采样色之间}的 CIEDE2000 色差均大于阈值：
\begin{equation}
\Delta E_{00}(c_i, c_j) \geq 3 \quad \forall\; i \neq j, \quad i, j \in \{1, \dots, 10\}
\end{equation}
该检验剔除了色差过小、感知上无法区分的退化色彩映射。

\subsection{双窗口均匀性检验}
通过质量检验的色彩映射必须同时满足\textbf{小窗口}和\textbf{大窗口}两个层级的均匀性约束。

从 256 色中均匀采样 $N_s = 30$ 个颜色，构成等间距序列 $\{c_1, c_2, \dots, c_{30}\}$。

\textbf{小窗口检验}（检测局部突变）：
\begin{equation}
\min_{|i - j| \geq k_S} \Delta E_{00}(c_i, c_j) \geq \tau_S \quad \text{其中 } k_S = 2,\; \tau_S = 3
\end{equation}
确保间隔 $\geq 2$ 的采样点之间有足够色差，防止"色带"——一段视觉上无变化的区域。

\textbf{大窗口检验}（检测全局单调性）：
\begin{equation}
\min_{|i - j| \geq k_L} \Delta E_{00}(c_i, c_j) \geq \tau_L \quad \text{其中 } k_L = 5,\; \tau_L = 10
\end{equation}
确保较大间隔的采样点之间有显著色差，保障色彩映射整体上能有效传达数值差异。


\section{优化策略：贪心搜索 + 局部精炼}

\subsection{问题形式化}
对于每个参数组合 $(H_{target}, C_{pattern}, L_{pattern})$，目标是找到最小的扰动偏移 $(\delta_C^*, \delta_L^*)$，使得存在至少一条通过双窗口质量检验的色彩映射：
\begin{equation}
(\delta_C^*, \delta_L^*) = \arg\min_{\delta_C, \delta_L \geq 0} \sqrt{\delta_C^2 + \delta_L^2}
\quad \text{s.t.} \quad \exists\; \text{colormap}(\delta_C, \delta_L) \text{ passes quality check}
\end{equation}

直觉：偏移越小，生成的色彩映射越接近"理想"控制点模式（色度/亮度值域更完整），感知表达能力越强。

\subsection{扰动偏移搜索空间}
系统预定义一组离散的 $(\delta_C, \delta_L)$ 候选对，按欧氏距离排序：
\begin{itemize}
    \item 细粒度区间：$\delta \in \{0, 2, 4, \dots, 20\}$，步长 $= 2$（$11 \times 11 = 121$ 组）
    \item 中粒度区间：$\delta \in \{25, 30, \dots, 50\}$，步长 $= 5$（$6 \times 6 = 36$ 组）
    \item 粗粒度区间：$\delta_C \in \{60, 70, \dots, 100\}$，$\delta_L \in \{60, 70\}$，步长 $= 10$（$5 \times 2 = 10$ 组）
\end{itemize}
所有候选对按 $\sqrt{\delta_C^2 + \delta_L^2}$ 升序排列，总计约 167 个偏移对。

\subsection{两阶段搜索算法}

\begin{algorithm}[H]
\caption{贪心搜索 + 局部精炼}
\begin{algorithmic}[1]
\State \textbf{Phase 1: 贪心搜索（找到第一个可行解）}
\For{$(\delta_C, \delta_L)$ in sorted offsets (distance ascending)}
    \State 调整控制点：$C' \leftarrow \text{adjust}(C_{pattern}, \delta_C)$, $L' \leftarrow \text{adjust}(L_{pattern}, \delta_L)$
    \State 调用网格搜索生成候选集
    \If{存在通过双窗口检验的候选}
        \State 记录当前 $(\delta_C, \delta_L)$ 为 $(\delta_C^{feas}, \delta_L^{feas})$
        \State \textbf{break}
    \EndIf
\EndFor
\State
\State \textbf{Phase 2: 局部精炼（在可行解附近搜索更优解）}
\State 定义搜索窗口 $W = 7$（QUALITY 模式）
\For{$dc = \max(0, \delta_C^{feas} - W)$ to $\min(100, \delta_C^{feas} + W)$}
    \For{$dl = \max(0, \delta_L^{feas} - W)$ to $\min(70, \delta_L^{feas} + W)$}
        \If{$\sqrt{dc^2 + dl^2} < \sqrt{(\delta_C^*)^2 + (\delta_L^*)^2}$} \Comment{仅尝试更优的偏移}
            \State 生成并检验
            \If{通过质量检验}
                \State 更新 $(\delta_C^*, \delta_L^*) \leftarrow (dc, dl)$
            \EndIf
        \EndIf
    \EndFor
\EndFor
\State \Return 最优偏移对应的色彩映射
\end{algorithmic}
\end{algorithm}

\subsection{色相重试机制}
如果在当前色相目标下，穷尽所有偏移对仍然找不到可行解，系统会从 $H_{target}$ 附近随机偏移重试：
\begin{equation}
H_{actual} = H_{target} + \mathcal{U}(-10, 10)
\end{equation}
最大重试 100 次。这处理了某些色相角度在 HCL-to-sRGB 映射中色域特别狭窄的情况。


\section{最优选择（后处理）}

当同一参数组合 $(H_{target}, C_{pattern}, L_{pattern})$ 产生多个候选时（例如多次运行），系统从中选出\textbf{最优}个体。选择标准为：

\begin{enumerate}
    \item \textbf{主排序}：扰动距离 $d = \sqrt{\delta_C^2 + \delta_L^2}$ 最小者优先（越接近理想控制点越好）
    \item \textbf{次排序}（平局时）：重试次数 $r$ 最少者优先（搜索越容易 $\Rightarrow$ 参数空间越稳定）
\end{enumerate}

最终输出按 $H_{target} \to L_{pattern} \to C_{pattern}$ 排序的 48 条色彩映射，ID 重新编号为 $0 \sim 47$。

\end{document}
